% ===================================================================
%  جدول نمبر 5 — گھر اور اُس کی تعمیر۔
%
%  Traduction, faite en 2026, en ourdou d'aujourd'hui, sur l'ido de
%  texto/io. Regles de la colonne : voir 10-jadval-01.tex. Le tableau
%  entier est un DIALOGUE et aucun alinea n'est numerote par
%  l'auteur : toutes les cles portent « 01 », le rang suivant la
%  replique, comme en ido. Ecrit avec outils/ordo.py sous les yeux.
%
%  LE TABLEAU LE PLUS LONG DU LIVRET — soixante-neuf blocs, plus de
%  cent quarante renvois —, et le cinquieme ecrit avec la liste de
%  parigi.py en main des la premiere ligne.
%
%  ONZE RETOURNEMENTS SUR TRENTE-SEPT PAIRES, comme au tamoul. Le
%  coreen en avait compte douze. C'est le taux le plus BAS de la
%  colonne jusqu'ici, et la raison est celle du marche telougou : un
%  chantier s'enumere. On nomme les corps de metier l'un apres
%  l'autre, et les outils de chacun a la suite.
%
%  LES METIERS DU BATIMENT SONT PERSANS ET INDIENS A PARTS EGALES, ET
%  LE PARTAGE N'EST PAS OU ON L'ATTEND. معمار le macon, نجار le
%  charpentier, لوہار le forgeron, ترکھان le menuisier, مالی le
%  jardinier, راج مزدور le manœuvre : deux mots arabes, deux mots
%  indiens, deux composes. Les outils, eux, sont indiens presque
%  tous — کرنی la truelle, ساہول le fil a plomb, رندہ le rabot, چھینی
%  le ciseau, نہائی l'enclume, سنڈسی les tenailles, ریتی la lime,
%  کنڈی le verrou, ٹھاٹھ l'echafaudage, کھپرا la tuile. La langue
%  savante nomme l'homme, la langue de la maison nomme la chose qu'il
%  tient.
%
%  ET L'OURDOU DISTINGUE CE QUE LE TAMOUL CONFONDAIT. Le tamoul
%  n'avait qu'un mot, தச்சர், pour le charpentier (115) et le
%  menuisier, et avait du qualifier. L'ourdou a نجار pour l'un et
%  ترکھان pour l'autre, deux mots entiers venus de deux cotes — l'un
%  arabe, l'autre indien. Le coreen distinguait de meme, 목수 et
%  소목장이. Trois colonnes, trois solutions, et une seule qui doive
%  qualifier faute de mot.
%
%  LE BALCON (28) EST LE SIXIEME PIEGE, ET IL SE RETOURNE COMME LE
%  VERANDO DU TABLEAU 4. Le tamoul avait ecrit l'emprunt பால்கனி
%  faute de mot juste ; l'ourdou a بالکونی — emprunt lui aussi — mais
%  aussi چھجا, l'auvent saillant de la maison indienne, qui n'est pas
%  la meme chose. La colonne ecrit بالکونی et le declare : ici
%  l'ourdou emprunte, la ou il pretait au tableau precedent. Une
%  langue n'est pas riche ou pauvre en bloc, elle l'est objet par
%  objet.
%
%  DEUX ECARTS AU PREMIER PASSAGE, ET CE SONT EXACTEMENT LES DEUX QUI
%  AVAIENT TROMPE LA COLONNE TAMOULE, AUX MEMES BLOCS. Au 01-30,
%  « indutas per gipso (70) la plafoni (26) » : l'ido met le moyen
%  avant l'objet, et la phrase naturelle veut l'inverse dans les deux
%  langues. Au 01-38, « la tapetisto (107) qua esas sur skalo (106),
%  sur la balkono (28), pozas storo (29) » : la relative s'intercale
%  entre le sujet et son verbe, et l'on est tente de la rejeter en
%  fin de phrase — ce qui remonte le storo.
%
%  Ni l'un ni l'autre n'est un rapport modificateur–tete : parigi.py
%  ne les liste pas et n'a pas a les lister. Que deux langues sans
%  parente butent sur les deux MEMES blocs d'un tableau de cent
%  quarante renvois n'est donc pas une coincidence de traducteur :
%  ce sont deux endroits ou l'ido ordonne autrement que ne le ferait
%  n'importe quelle phrase a tete finale. ordo.py les montre ; il
%  reste a les lire.
%
%  LE FAC-SIMILE REEMPLOIE TROIS DE SES RENVOIS et la source ne bouge
%  pas : (107) est le porteur de mortier au bloc 01-21 et le tapissier
%  au bloc 01-38 ; (33) est le verando au bloc 01-30 et les mansardes
%  au bloc 01-55 ; (1) est la facade ET la salle de bain dans la meme
%  phrase du bloc 01-53.
% ===================================================================

%%K t05-apar-1 apar
\VUsaut{6.0mm}
\VUtitre{15.0pt}{60}{جدول نمبر 5}
\VUsaut{1.4mm}
\VUtitre{11.4pt}{0}{گھر اور اُس کی تعمیر۔}
\VUsaut{2.2mm}

%%K t05-tit-1 sub
\VUpk{10.2pt}{40}{اچھی ملاقات۔}
\VUsaut{3.0mm}

%%K t05-01-1 p
\textsc{یوآنیس}۔ --- ماموں جان، میں بہت جاننا چاہتا ہوں کہ
\VUgras{گھر} کیسے بنایا جاتا ہے۔

%%K t05-01-2 p
\textsc{ماموں} \textsuperscript{(80)}۔ --- یہ تو بہت آسان ہے، میرے
دوست ؛ میرے ساتھ آؤ۔ میں تمھیں وہ گھر دکھاتا ہوں جو برنارڈس صاحب
\textsuperscript{(79)} بنوا رہے ہیں اور جس میں میں رہوں گا۔

%%K t05-01-3 p
\textsc{یوآنیس}۔ --- اِس گھر کا \VUgras{نقشہ} \textsuperscript{(76)}
کس نے بنایا؟

%%K t05-01-4 p
\textsc{ماموں}۔ --- \VUgras{ماہرِ تعمیر} \textsuperscript{(75)} نے۔
اُنھوں نے بہت صاف خاکہ پیش کیا، جو منظور ہو گیا ؛ پھر
\VUgras{ٹھیکیدار} \textsuperscript{(77)} نے کام شروع کیا۔

%%K t05-01-5 p
\textsc{یوآنیس}۔ --- ٹھیکیدار کون ہیں؟

%%K t05-01-6 p
\textsc{ماموں}۔ --- وہ بلانکو صاحب ہیں، تمھارے دوست یاکوبس کے والد۔
وہ ماہرِ تعمیر روفو صاحب کے ساتھ مل کر مزدوروں کو چلاتے ہیں۔

%%K t05-01-7 p
لو! ٹھیک وقت پر بلانکو صاحب اپنا \VUgras{پیمانہ}
\textsuperscript{(78)} لیے آ رہے ہیں ؛ اور روفو صاحب بھی اپنا نقشہ لیے
آ رہے ہیں۔

%%K t05-01-8 p
السلام علیکم، صاحبو۔ مزاج کیسے ہیں؟

%%K t05-01-9 p
\textsc{روفو صاحب}۔ --- بہت اچھے، شکریہ۔ السلام علیکم، یوآنیس ؛ یہ
بچہ کتنا بڑا ہو گیا!

%%K t05-01-10 p
\textsc{ماموں}۔ --- ہاں، واقعی، یہ چھوٹا آدمی ہے۔ بہت سنجیدہ ہے ؛ جو
کچھ دیکھتا ہے، اُس کے بارے میں اِسے بتانا پڑتا ہے۔ آج یہ جاننا چاہتا
ہے کہ گھر کیسے بنایا جاتا ہے۔

%%K t05-tit-2 sub
\VUpk{10.2pt}{40}{مزدور۔}
\VUsaut{3.0mm}

%%K t05-01-11 p
\textsc{بلانکو صاحب}۔ --- آؤ میرے ساتھ، ننھے دوست ؛ ابھی تمھیں سمجھا
دیتا ہوں، کیونکہ مجھے وہ لڑکے بہت پسند ہیں جو سیکھنا چاہتے ہیں۔

%%K t05-01-12 p
تم وہ مزدور دیکھ رہے ہو جس کے ہاتھ میں \VUgras{بیلچہ}
\textsuperscript{(82)} ہے : وہ \VUgras{مٹی کا مزدور}
\textsuperscript{(81)} ہے۔ وہ پانی کے پائپ ڈالنے کے لیے
\VUgras{خندق} \textsuperscript{(46)} کھود رہا ہے۔ اُس نے اُس جگہ کی
زمین بھی کھودی جہاں گھر ہے، تاکہ معمار چار \VUgras{دیواریں} اٹھا
سکے۔ اُن دیواروں کا جو حصہ زمین کے اندر ہے، اُسے \VUgras{بنیاد}
\textsuperscript{(2)} کہتے ہیں۔ بنیادوں کے بیچ کی جگہ
\VUgras{تہ خانہ} \textsuperscript{(3)} بناتی ہے۔ اُس تہ خانے میں
\VUgras{تہ خانے کی کھڑکی} \textsuperscript{(5)} نامی ایک چھوٹی کھڑکی،
ایک \VUgras{دروازہ} \textsuperscript{(4)} اور ایک زینہ ہے۔

%%K t05-01-13 p
\VUgras{معمار} \textsuperscript{(108)} نے \VUgras{دروازوں}
\textsuperscript{(8)} اور \VUgras{کھڑکیوں} \textsuperscript{(17)} کے
لیے سوراخ چھوڑتے ہوئے دیواریں بناتا رہا۔

%%K t05-01-14 p
\textsc{یوآنیس}۔ --- مگر وہ معمار مجھے نظر نہیں آ رہا!

%%K t05-01-15 p
\textsc{بلانکو صاحب}۔ --- اب اُس نے گھر میں کام ختم کر لیا۔ وہ
\VUgras{اصطبل} \textsuperscript{(54)} کے قریب، اپنے \VUgras{ٹھاٹھ}
\textsuperscript{(110)} پر ہے ؛ وہ \VUgras{دھوبی خانے}
\textsuperscript{(56)} کی دیوار بنا رہا ہے۔

%%K t05-01-16 p
اُس کے پاس کرنی، تسلا اور \VUgras{ساہول} \textsuperscript{(109)} ہیں۔
مگر یہ نام تمھیں یاد نہیں رہیں گے۔

%%K t05-01-17 p
\textsc{یوآنیس}۔ --- ضرور رہیں گے ؛ کیونکہ ابھی اپنے ماموں سے چھوٹی
کرنی اور تسلا مانگوں گا ؛ اور ساہول خود ایک ڈوری سے بنا لوں گا۔

%%K t05-tit-3 sub
\VUsaut{3.0mm}
\VUpk{10.2pt}{40}{تعمیر کا سامان۔}
\VUsaut{3.0mm}

%%K t05-01-18 p
\textsc{ماموں}۔ --- آہا! بہت خوب۔ مگر بتاؤ یوآنیس، گھر بنانے کے لیے
کیا چاہیے؟

%%K t05-01-19 p
\textsc{یوآنیس}۔ --- \VUgras{تراشے ہوئے پتھر} \textsuperscript{(59)}
اور \VUgras{گارا} \textsuperscript{(65)} چاہیے۔

%%K t05-01-20 p
\textsc{ماموں}۔ --- ٹھیک۔ بڑی ٹوپی پہنے اپنی \VUgras{گاڑی}
\textsuperscript{(136)} پر بیٹھے اُس آدمی کو دیکھو۔ وہ
\VUgras{گاڑی بان} \textsuperscript{(135)} ہے۔ روز یہاں \VUgras{پتھر
کے بلاک} \textsuperscript{(60)} لاتا ہے۔ صحن میں اپنی گاڑی خالی کرتا
ہے ؛ \VUgras{پتھر تراشنے والے} \textsuperscript{(83)} اُن بلاکوں کو
تراشتے ہیں اور اپنے \VUgras{پتھر کے آرے} \textsuperscript{(85)} سے
چیرتے ہیں۔ \VUgras{راج مزدور} \textsuperscript{(146)} اُنھیں معمار کے
پاس لے جاتا ہے ؛ معمار اُنھیں جگہ پر رکھتا ہے۔

%%K t05-01-21 p
اُسی وقت \VUgras{گارا گھولنے والا} \textsuperscript{(87)} گارا تیار
کرتا ہے۔ اُسے \VUgras{ریت} \textsuperscript{(62)}، \VUgras{چونا}
\textsuperscript{(63)} اور \VUgras{پانی} \textsuperscript{(64)} چاہیے۔
چونا اُس کے پاس ایک \VUgras{بورے} \textsuperscript{(71)} میں ہے ؛
\VUgras{معمار کا مددگار} \textsuperscript{(111)} ایک \VUgras{ٹھیلے}
\textsuperscript{(112)} میں ریت لاتا ہے ؛ \VUgras{شاگرد}
\textsuperscript{(113)} نلکے (58) کے پاس ہے۔ وہ ایک \VUgras{بالٹی}
\textsuperscript{(114)} بھر رہا ہے۔ گارا جلد تیار ہو جاتا ہے۔
\VUgras{گارا ڈھونے والا} \textsuperscript{(107)} اُسے لے کر سیڑھی کے
راستے ٹھاٹھ پر چڑھاتا ہے ؛ وہاں گھر کا وہ پہلو ختم کرنے والا معمار
کام کر رہا ہے۔

%%K t05-01-22 p
اور اب بتاؤ، \VUgras{چھت} \textsuperscript{(31)} پر کام کرنے والے
مزدور کو کیا کہتے ہیں؟

%%K t05-01-23 p
\textsc{یوآنیس}۔ --- وہ \VUgras{چھت بنانے والا} \textsuperscript{(118)}
ہے۔

%%K t05-tit-4 sub
\VUsaut{3.0mm}
\VUpk{10.2pt}{40}{گھر کی چھت۔}
\VUsaut{3.0mm}

%%K t05-01-24 p
\textsc{بلانکو صاحب}۔ --- ہاں ؛ مگر پہلے \VUgras{نجار}
\textsuperscript{(115)} آتا ہے۔

%%K t05-01-25 p
نجار \VUgras{لکڑی کے ڈھانچے} \textsuperscript{(35)} پر کام کرتا ہے۔ وہ
\VUgras{شہتیروں} \textsuperscript{(37)} کو \VUgras{لکڑی کی میخوں}
\textsuperscript{(117)} سے جوڑتا ہے۔ وہاں اوپر، چوڑے مخملی پاجامے
پہنے وہ مزدور نجار ہیں۔ ایک چھوٹا شہتیر پکڑے ہوئے ہے ؛ دوسرا اپنے
\VUgras{کلہاڑے} \textsuperscript{(116)} سے میخ کاٹ رہا ہے۔
\VUgras{دھوئیں کی چمنی} \textsuperscript{(41)} کے قریب جو ہے، وہ چھت
بنانے والا ہے۔

%%K t05-noto-1 noto
\VUnotes{143.2mm}{%
(*) \VUgras{Balk-o} = ج۔ \textit{Balken}، ا۔ \textit{joist}، ف۔
\textit{solive}، ا۔ \textit{travicello}، ر۔ \textit{balka}، ہ۔ پ۔
\textit{viga}۔
یہ فنی مادہ اب تک قبول نہیں کیا گیا ؛ مگر ف۔ \textit{solive} کے
معنی کے ترجمے کے لیے تجویز کے قابل ہے۔ وہ trabo (ف۔
\textit{poutre}) نہیں، نہ trabeto۔ وہ چوکور تراشی ہوئی لکڑی کی
شہتیری ہے جو فرش اور چھت کو تھامتی ہے۔%
}

%%K t05-01-25 p suite
وہ (*) شہتیریوں پر پٹیاں اور پٹیوں پر \VUgras{سلیٹ کے تختے}
\textsuperscript{(66)} جماتا ہے۔ کبھی کبھی کھپرے ڈالتا ہے۔

%%K t05-01-26 p
اصطبل کی چھت \VUgras{کھپروں سے ڈھکی} \textsuperscript{(55)} ہے ؛ گھر
کی چھت \VUgras{سلیٹ سے ڈھکی} \textsuperscript{(32)} ہے ؛ برآمدے کی چھت
\VUgras{جست} \textsuperscript{(24)} کی ہے۔

%%K t05-01-27 p
\textsc{یوآنیس}۔ --- کیا جست کی چھتوں پر بھی چھت بنانے والا ہی کام
کرتا ہے؟

%%K t05-01-28 p
\textsc{بلانکو صاحب}۔ --- نہیں ؛ \VUgras{جست کا کاریگر}
\textsuperscript{(121)} یا \VUgras{سیسے کا کاریگر} کرتا ہے ؛ وہی گھر
کی چھت پر \VUgras{روشن دان} \textsuperscript{(38)} لگاتا ہے۔
\VUgras{برساتی نالیاں} \textsuperscript{(43)} اور \VUgras{چھت کے
پرنالے} \textsuperscript{(42)} بھی وہی لگاتا ہے۔ جست اور ٹین کو وہ
ٹانکا لگاتا ہے۔ \VUgras{آسمانی بجلی کا محافظ} \textsuperscript{(40)}
اور \VUgras{ہوا کا رخ بتانے والا} \textsuperscript{(39)} اُسی نے
لگائے ؛ وہ \VUgras{کھپریل کی نوکدار دیوار} \textsuperscript{(36)} پر
ہیں۔

%%K t05-01-29 p
\textsc{یوآنیس}۔ --- تو کیا گھر پورا ہو گیا؟

%%K t05-tit-5 sub
\VUsaut{3.0mm}
\VUpk{10.2pt}{40}{اوزار۔}
\VUsaut{3.0mm}

%%K t05-01-30 p
\textsc{بلانکو صاحب}۔ --- بالکل نہیں۔ \VUgras{پلستر کرنے والا}
\textsuperscript{(99)} \VUgras{اینٹوں} \textsuperscript{(61)} سے وہ
دیواریں بناتا ہے جو گھر کو کئی کمروں میں بانٹتی ہیں۔ وہ
\VUgras{پلستر} \textsuperscript{(70)} سے \VUgras{چھتوں}
\textsuperscript{(26)} اور دیواروں کو لیپتا ہے۔ اب وہ \VUgras{برآمدے}
\textsuperscript{(33)} کی چھت پر کام کر رہا ہے۔ معمار کی طرح اُس کے
پاس بھی \VUgras{کرنی} \textsuperscript{(100)} اور \VUgras{تسلا}
\textsuperscript{(101)} ہیں۔

%%K t05-01-31 p
پھر \VUgras{ترکھان} دروازوں، کھڑکیوں، \VUgras{لکڑی کے فرشوں}
\textsuperscript{(16)} اور \VUgras{کھڑکیوں کے پٹوں}
\textsuperscript{(19)} پر کام کرتا ہے۔

%%K t05-01-32 p
اب وہ صحن میں اپنے \VUgras{کام کے تختے} \textsuperscript{(90)} پر کام
کر رہا ہے۔ اُس کے پاس \VUgras{آرا} \textsuperscript{(94)} ہے — وہ
\VUgras{لکڑی} (69) چیرتا ہے —، اور \VUgras{رندہ}
\textsuperscript{(91)}، \VUgras{شکنجہ} \textsuperscript{(92)}،
\VUgras{چھینی} \textsuperscript{(94 bis)}، \VUgras{پرکار}
\textsuperscript{(95)}، لکڑی کا \VUgras{ہتھوڑا} \textsuperscript{(95 bis)}
ہیں۔

%%K t05-01-33 p
\textsc{یوآنیس}۔ --- آہا! مگر معماری سے ترکھانی زیادہ مزے کی ہے۔
ماموں جان، مجھے کام کا تختہ، آرا اور رندہ خرید دیں گے؟ میں جلد
میدور کے لیے ایک جھونپڑی بناؤں گا۔

%%K t05-01-34 p
\textsc{ماموں}۔ --- ہاں، بیٹا۔

%%K t05-01-35 p
\textsc{یوآنیس}۔ --- میں کتنا خوش ہوں! اور داخلی دروازے کے پاس جو
کھڑا ہے، وہ کون ہے؟

%%K t05-tit-6 sub
\VUsaut{3.0mm}
\VUpk{10.2pt}{40}{رنگ کرنا۔}
\VUsaut{3.0mm}

%%K t05-01-36 p
\textsc{بلانکو صاحب}۔ --- وہ \VUgras{رنگ ساز} \textsuperscript{(102)}
ہے۔ \VUgras{رنگ} \textsuperscript{(103)} کی ہانڈی اور اپنا
\VUgras{برش} \textsuperscript{(104)} لیے اپنی سیڑھی پر کھڑا ہے۔ داخلی
دروازے کی لکڑی زیادہ دیر چلے، اِس لیے اُس پر رنگ کر رہا ہے۔

%%K t05-01-37 p
کمروں میں کاغذ چپکانا بھی وہی کرتا ہے۔

%%K t05-01-38 p
\VUgras{کاغذ چپکانے والا} \textsuperscript{(107)} — وہ ایک
\VUgras{سیڑھی} \textsuperscript{(106)} پر، \VUgras{بالکونی}
\textsuperscript{(28)} میں ہے — \VUgras{پردہ} \textsuperscript{(29)}
لگا رہا ہے۔

%%K t05-01-39 p
بڑی کھڑکی کے پاس کھڑا مزدور \VUgras{شیشہ ساز} \textsuperscript{(96)}
ہے۔ وہ \VUgras{شیشے کے تختے} \textsuperscript{(18)} \VUgras{لیس}
\textsuperscript{(73)} سے جماتا ہے۔ تہ خانے کے دروازے کے پاس گھٹنے
ٹیکے وہ دوسرا مزدور \VUgras{قفل ساز} \textsuperscript{(97)} ہے۔ وہ
\VUgras{قفل} \textsuperscript{(6)} لگا رہا ہے۔ اُس کی \VUgras{ریتی}
\textsuperscript{(98)} اُس کے پاس ہے۔

%%K t05-01-40 p
\textsc{یوآنیس}۔ --- اپنے \VUgras{ہتھوڑے} \textsuperscript{(84)} سے
لوہے پر مارنے والا وہ مزدور بھی قفل ساز ہے؟

%%K t05-01-41 p
\textsc{بلانکو صاحب}۔ --- زیادہ ٹھیک کہوں تو وہ لوہار (125) ہے۔ وہ
\VUgras{لوہے کا سامان} \textsuperscript{(67)} گھڑتا ہے ؛ وہ سامان
\VUgras{بجلی والے} \textsuperscript{(124)} کے لیے ہے۔ وہ بجلی والا
\VUgras{گاڑی خانے} \textsuperscript{(51)} کی چھت پر ہے۔ لوہار کے پاس
\VUgras{بھٹی} \textsuperscript{(126)}، \VUgras{نہائی}
\textsuperscript{(127)} اور \VUgras{سنڈسی} \textsuperscript{(128)}
ہیں۔

%%K t05-01-42 p
بجلی والا ٹیلیفون کا \VUgras{تانبے} \textsuperscript{(44)} کا
\VUgras{تار} لگا رہا ہے۔

%%K t05-tit-7 sub
\VUpk{10.2pt}{40}{باغبانی۔}
\VUsaut{3.0mm}

%%K t05-01-43 p
\textsc{ماموں}۔ --- \VUgras{صحن} \textsuperscript{(45)} کے پیچھے،
\VUgras{لوہے کے جنگلے} \textsuperscript{(47)} اور \VUgras{بڑے پھاٹک}
\textsuperscript{(48)} کے قریب تمھیں \VUgras{مالی}
\textsuperscript{(129)} نظر آتا ہے۔ وہ \VUgras{چھوٹا باغ}
\textsuperscript{(49)} تیار کر رہا ہے۔ اپنے \VUgras{بیلچے}
\textsuperscript{(131)} سے مٹی کھودی، بیج بوتا ہے، \VUgras{لوہے کے
پنجے} \textsuperscript{(130)} سے ہموار کرتا ہے۔ پھر اپنے
\VUgras{فوارے} \textsuperscript{(132)} سے \VUgras{کیاریوں}
\textsuperscript{(150)} کو پانی دے گا، اور پودے بڑھیں گے۔

%%K t05-01-44 p
گھوڑے کی دیکھ بھال کر کے اُسے چارا دے چکنے والا \VUgras{سائیس}
\textsuperscript{(133)} \VUgras{گاڑی} \textsuperscript{(52)} کو
اسفنج، \VUgras{ہاتھ کے پمپ} \textsuperscript{(134)} اور ایک بالٹی
پانی سے دھو رہا ہے۔ پھر اُسے گاڑی خانے میں واپس لے جائے گا اور موٹی
\VUgras{کنڈی} \textsuperscript{(53)} سے دروازہ بند کرے گا۔

%%K t05-01-45 p
لو، میرے خیال میں تم خوش ہو ؛ تمھیں کافی وضاحتیں مل گئیں۔

%%K t05-01-46 p
\textsc{یوآنیس}۔ --- ہاں ماموں جان ؛ مگر گھر کے اندر بھی شوق سے
دیکھوں گا۔

%%K t05-01-47 p
\textsc{ماموں}۔ --- وہ کل ہو گا۔ روفو صاحب تمھیں نقشہ سمجھائیں گے اور
ہر کمرے کا نام بتائیں گے۔

%%K t05-tit-8 sub
\VUsaut{3.0mm}
\VUpk{10.2pt}{40}{گھر کی اندرونی ترتیب۔}
\VUsaut{2.4mm}

%%K t05-apar-2 apar
{\centering\textit{(نقشہ دیکھیے۔)}\par}
\VUsaut{2.4mm}

%%K t05-01-48 p
\textsc{ماہرِ تعمیر}۔ --- آؤ یوآنیس، ابھی تمھیں گھر کے مختلف حصے
گھما دیتا ہوں۔

%%K t05-01-49 p
\VUgras{زمینی منزل} \textsuperscript{(7)} میں داخل ہوں۔ یہ رہا
\VUgras{گھنٹی} \textsuperscript{(11)} کا بٹن۔ ہم تین \VUgras{زینوں}
\textsuperscript{(14)} پر چڑھتے ہیں ؛ وہ \VUgras{داخلی زینے}
\textsuperscript{(9)} کے ہیں۔ پھر \VUgras{دہلیز} \textsuperscript{(10)}
پار کرتے ہیں ؛ وہ \VUgras{داخلی دروازے} \textsuperscript{(8)} کی ہے۔
اب ہم \VUgras{زینے} \textsuperscript{(13)} کے نیچے ہیں۔

%%K t05-01-50 p
\textsc{یوآنیس}۔ --- یہ راہداری ہے۔

%%K t05-01-51 p
\textsc{ماہرِ تعمیر}۔ --- نہیں، یہ \VUgras{ڈیوڑھی}
\textsuperscript{(12)} ہے۔ \VUgras{راہداری} \textsuperscript{(a)} سرے
پر ہے۔ دائیں طرف یہ رہے \VUgras{بیٹھک} \textsuperscript{(b)} اور
\VUgras{کھانے کا کمرہ} \textsuperscript{(c)} ؛ کھانے کے کمرے کے سامنے
\VUgras{باورچی خانہ} \textsuperscript{(d)} اور \VUgras{خدمت گاہ}
\textsuperscript{(e)} ؛ پھر آگے \VUgras{لکھنے کا کمرہ}
\textsuperscript{(f)}۔ \VUgras{برآمدہ} \textsuperscript{(23)} بیٹھک کے
قریب ہے ؛ اُس کا اپنا \VUgras{شیشے کا دروازہ} \textsuperscript{(25)}
ہے۔

%%K t05-01-52 p
اب زینے پر چڑھتے ہیں۔ \VUgras{جنگلے} \textsuperscript{(15)} کو تھام
لو اور زینے گنو۔ وہ چودہ ہیں۔ یہ رہا \VUgras{زینے کا چبوترہ}
\textsuperscript{(m)} ؛ ہم \VUgras{پہلی منزل} \textsuperscript{(27)} پر
ہیں۔

%%K t05-tit-9 sub
\VUsaut{3.0mm}
\VUpk{10.2pt}{40}{پہلی منزل۔}
\VUsaut{3.0mm}

%%K t05-01-53 p
زینے کے سامنے \VUgras{بیت الخلا} \textsuperscript{(20)} اور
\VUgras{بناؤ سنگھار کا کمرہ} \textsuperscript{(j)} ہیں۔ دائیں طرف
خوبصورت \VUgras{سونے کا کمرہ} \textsuperscript{(g)} ؛ اُس کی گھر کے
\VUgras{اگلے رخ} \textsuperscript{(1)} پر دو کھڑکیاں ہیں۔ بائیں طرف
\VUgras{غسل خانہ} \textsuperscript{(1)} اور \VUgras{مہمانوں کا کمرہ}
\textsuperscript{(i)} ؛ اُس کمرے میں بڑا \VUgras{طاق نما}
\textsuperscript{(h)} ہے۔ سونے کے کمرے کے قریب \VUgras{بچوں کا کمرہ}
\textsuperscript{(k)}۔ وہ ایک کشادہ \VUgras{چھت} \textsuperscript{(21)}
سے لگا ہے۔ اُس چھت کے نیچے ایک اور بیت الخلا (20) ہے ؛ اور دیوار پر
یہ رہی وہ \VUgras{گھنٹی} \textsuperscript{(22)} جو رات کے کھانے کے لیے
بجتی ہے۔

%%K t05-01-54 p
\textsc{یوآنیس}۔ --- \VUgras{دوسری منزل} پر چلیں۔

%%K t05-01-55 p
\textsc{ماہرِ تعمیر}۔ --- دوسری منزل نہیں ہے۔ چھت کے نیچے صرف
\VUgras{چھجے کے کمرے} \textsuperscript{(33)} اور اناج کا گودام (34)
ہیں۔ ہم نے دیکھ لیا، اب اتر سکتے ہیں۔

%%K t05-01-56 p
\textsc{یوآنیس}۔ --- شکریہ صاحب، یہ بہت دلچسپ تھا۔ جو کچھ دیکھا، وہ
سب ابھی اپنے والد کو سناؤں گا۔

\VUsaut{4.0mm}
\VUfilet{22mm}
